\documentclass[10pt,a4paper]{amsart}
\usepackage[margin=19mm]{geometry}
\usepackage{amsmath,amssymb,mathtools}
\usepackage[scr=rsfs,cal=euler]{mathalfa}
\usepackage{xfrac}
\usepackage[unicode,hidelinks,bookmarks=false]{hyperref}
\newcommand{\ie}{i.e.\ }
\newcommand{\cf}{cf.\ }
\newcommand{\resp}{respectively\ }
\newcommand{\Subsection}{\subsection}
\providecommand{\nobreakdash}{\nobreak\hbox{-}}
\setlength{\emergencystretch}{2em}
\multlinegap0pt
\usepackage{bm}
\usepackage{bbm}
\usepackage{dsfont}
\usepackage{xspace}
\usepackage{cancel}
\usepackage{mathtools}
\usepackage{amsthm}
\makeatletter
\@ifundefined{thm}{\newtheorem{thm}{Theorem}}{}
\makeatother
\title[Multiple Riemann wave solutions of the general form of quasi]{Multiple Riemann wave solutions of the general form of quasilinear hyperbolic systems}
\author{A.M. Grundland, J. de Lucas}
\date{Source: arXiv:2203.15122v1 (CC BY 4.0)}
\begin{document}
\maketitle

\noindent\emph{Source and licence.} Multiple Riemann wave solutions of the general form of quasilinear hyperbolic systems. Version arXiv:2203.15122v1 (CC BY 4.0), distributed under the Creative Commons Attribution 4.0 International licence, as recorded in the arXiv licence metadata for this exact version and shown on the abstract page https://arxiv.org/abs/2203.15122v1. Full rights notice: licenses/arxiv-2203.15122-CC-BY-4.0.txt.\par\smallskip

\noindent\textit{Editorial scope.} Counted unit: the Thm whose source label is \texttt{None} in the article (source0.tex lines 1644--1651, source label \texttt{None}), delivered together with the article's own complete original proof of it (source0.tex lines 1653--1717, one \texttt{proof} environment of 65 lines). No mathematics is written, repaired or re-derived: statement and proof are the article's own text line for line. Required context retained uncounted: none. Every definition, display and label of those lines is retained unchanged. Omitted from this package and not redistributed: 1--1643, 1652--1652, 1718--2034. Nothing inside a retained excerpt is omitted or altered: every retained body is delivered exactly as the article's lines are written, so no omission receipt of any kind accompanies this package. Citations: the retained excerpts carry no citation command at all, so no bibliography entry of the article is transcribed and no reference is redistributed. Reference adaptation: none was needed and none was made; every reference inside the delivered unit resolves inside it, so no unresolved reference or citation is produced. Printed slips kept literal: no printed word, symbol, hypothesis or number of the counted statement or of its proof was altered. Layout and class: the article's own class is \texttt{birkjour} with packages including hyperref, setspace, amssymb, amsthm, amsmath, psfig, xcolor, enumerate, array, appendix, xy; the wrapper is the lane's pinned \texttt{amsart} editorial wrapper at 10pt on A4, which restates the theorem-like environments the retained text needs and adds the article's own macro definitions that the retained text needs (none), with extra wrapper packages (bm, bbm, dsfont, xspace, cancel, mathtools). The delivered numbering is the unit's own. Assets: no retained span contains a figure environment, a graphics-inclusion command, a PDF-page inclusion, a table or a TikZ picture; no image file is copied into this package, the package declares no asset dependency and no image file is redistributed with it. Licence: version arXiv:2203.15122v1 of this article is distributed under the Creative Commons Attribution 4.0 International licence, as recorded in the arXiv licence metadata for this exact version and shown on the abstract page https://arxiv.org/abs/2203.15122v1. A full rights notice is delivered at licenses/arxiv-2203.15122-CC-BY-4.0.txt. Characters: every retained excerpt is pure ASCII, so the delivered engine, XeLaTeX with its default Latin Modern text font, reports no missing character.
\begin{thm} {\bf (Modified Frobenius Theorem by rescaling)} Consider a family $X_1,\ldots,X_r$ of  vector fields on a $q$-dimensional manifold $N$ such that
	\begin{enumerate}
		\item  $X_1\wedge\ldots\wedge X_r\neq0$
		\item  $[X_i,X_j]=h_{i,j}^iX_i+h_{i,j}^jX_j$ 
	\end{enumerate} 
	for certain functions $h_{i,j}^k\in C^\infty(N)$ with $i,j,k=1,\ldots,r$. Then, there exist functions $f_1,\ldots,f_r\in C^\infty(N)$ such that $[f_iX_i,f_jX_j]=0$ for all $i,j=1,\ldots,r$. 
	
\end{thm}

\begin{proof}
	We proceed by induction. The result is immediate for $r=2$. Assume that our theorem is true for $r<q$ vector fields on $\mathbb{R}^q$ and let us prove that our theorem remains true for $r+1$ vector fields. Let   $X_1,\ldots,X_r,X_{r+1}$ be vector fields on $\mathbb{R}^{q}$ satisfying the conditions 1 and 2. The first $r$ vector fields can be rescaled, by the induction hypothesis, in order to obtain $Y_1,\ldots,Y_r$ so  that $[Y_i,Y_j]=0$ with $i,j=1,\ldots,r$. Then, there exist coordinates $\chi_1,\ldots,\chi_r,\eta_1,\ldots,\eta_s$ on $\mathbb{R}^q$ such that
	$$
	X_i=\frac{\partial}{\partial \chi_i},\qquad i=1,\ldots,r.
	$$
	Let us determine whether  there exists a function $f_{r+1}\in C^\infty(N)$ such that 
	\begin{equation}\label{Eq:FC}
		[X_i,f_{r+1}X_{r+1}]\wedge X_{i}=0,\qquad i=1,\ldots,r.
	\end{equation}
	Condition 2 and (\ref{Eq:FC}) imply that
	$$
	[X_i,f_{r+1}X_{r+1}]=(X_if_{r+1})X_{r+1}+f_{r+1}(h_{i,r+1}^{r+1}X_{r+1}+h_{i,r+1}^iX_i),\,\, i=1,\ldots,r,
	$$
	for certain functions $h_{i,r+1}^{r+1},h_{i,r+1}^i$ on $N$ with $i=1,\ldots,r$. Hence, to satisfy the conditions (\ref{Eq:FC}), the function $f_{r+1}$ must be such that 
	\begin{equation}\label{Eq:SC}
		X_i\ln|f_{r+1}|=-h_{i,r+1}^{r+1},\qquad i=1,\ldots,r.
	\end{equation}
	Let us prove that the above system of partial differential equations admits a solution. Its compatibility condition reads
	$$
	\frac{\partial h_{i,r+1}^{r+1}}{\partial \chi_j}=\frac{\partial h_{j,r+1}^{r+1}}{\partial \chi_i},\qquad i,j=1,\ldots,r.
	$$
	To show that the above compatibility condition holds, let us consider the Jacobi identities
	$$
	[X_i,[X_j,X_{r+1}]]+[X_j,[X_{r+1},X_{i}]]+[X_{r+1},[X_i,X_j]]=0,\qquad i,j=1,\ldots,r.
	$$
	Since $[X_i,X_j]=0$ by the induction hypothesis, the Jacobi identities reduces to
	$$
	[X_i,[X_j,X_{r+1}]]-[X_j,[X_{i},X_{r+1}]]=0,\qquad i,j=1,\ldots,r.
	$$
	Using assumption 2, we obtain that 
	$$
	[X_i,[X_j,X_{r+1}]]=\frac{\partial h_{j,r+1}^{r+1}}{\partial \chi_i}X_{r+1}+h_{j,r+1}^{r+1}[X_i,X_{r+1}]+\frac{\partial h^{j}_{j,r+1}}{\partial \chi_i}X_j
	$$
	for any $i,j=1,\ldots,r$, and, moreover,
	\begin{multline*}
		0=[X_j,[X_i,X_{r+1}]]-[X_i,[X_j,X_{r+1}]]=\left(\frac{\partial h_{i,r+1}^{r+1}}{\partial \chi_j}-\frac{\partial h_{j,r+1}^{r+1}}{\partial \chi_i}\right)X_{r+1}\\+h_{i,r+1}^{r+1}[X_j,X_{r+1}]+\frac{\partial h^i_{i,r+1}}{\partial \chi_j}X_i-h_{j,r+1}^{r+1}[X_i,X_{r+1}]-\frac{\partial h^j_{j,r+1}}{\partial \chi_i}X_j\\=
		\left(\frac{\partial h_{i,r+1}^{r+1}}{\partial \chi_j}-\frac{\partial h_{j,r+1}^{r+1}}{\partial \chi_i}\right)X_{r+1}-\left(h_{j,r+1}^{r+1}h_{i,r+1}^{i}-\frac{\partial h^i_{i,r+1}}{\partial \chi_j}\right)X_i\\+\left(h^{r+1}_{i,r+1}h_{j,r+1}^j-\frac{\partial h^j_{j,r+1}}{\partial \chi_i}\right)X_j.
	\end{multline*}
	Since $X_1\wedge\ldots\wedge X_{r+1}\neq 0$, it follows that
	\begin{equation}\label{eq:ConS}
		\frac{\partial h_{j,r+1}^{r+1}}{\partial \chi_i}=\frac{\partial h_{i,r+1}^{r+1}}{\partial \chi_j},\qquad i,j=1,\ldots,r,\qquad i\neq j.
	\end{equation}
	These equalities imply that the system (\ref{Eq:SC}) has a solution  $f_{r+1}$ and therefore, $X_{r+1}$ can be rescaled so that $\widetilde{X}_{r+1}=f_{r+1}X_{r+1}$ satisfies
	\begin{equation}\label{Eq:TC}
		[X_i,\widetilde{X}_{r+1}]\wedge X_{i}=0,\qquad i=1,\ldots,r.
	\end{equation}
	Let us show that it is possible to rescale $X_1,\ldots, X_r$ so that they still commute among themselves and, in addition,  commute with $X_{r+1}$. In fact,  $[f_iX_i,f_jX_j]=0$ for every $i,j=1,\ldots,r$ if and only if 
	$$
	f_i=f_i(\chi_i,\eta),\qquad i=1,\ldots,r.
	$$
	The conditions (\ref{Eq:TC}) for $\widetilde{X}_{r+1}$ imply that the latter takes the form
	$$
	\widetilde{X}_{r+1}=\sum_{i=1}^rf^i_{r+1}(\chi_i,\eta_1,\ldots,\eta_s)\frac{\partial}{\partial \chi_i}+\sum_{\mu=1}^sf^\mu_{r+1}(\eta_1,\ldots,\eta_s)\frac{\partial}{\partial \eta_\mu}
	$$
	for some functions of the form $f_{r+1}^i(\chi_i,\eta_1,\ldots,\eta_s)$ and $f^\mu_{r+1}(\eta_1,\ldots,\eta_s)$ for $i=1,\ldots,r$ and   $\mu=1,\ldots,s$. 
	If $[f_iX_i,\widetilde{X}_{r+1}]=0$, the function $f_i$ must satisfy the partial differential equation
	\begin{equation}\label{Eq:PDEFro}
		0=\widetilde {X}_{r+1}f_i-f_i\frac{\partial f^i_{r+1}}{\partial \chi_i}=f^i_{r+1}(\chi_i,\eta)\frac{\partial f_i}{\partial \chi_i}+\sum_{\mu=1}^sf^\mu_{r+1}(\eta)\frac{\partial f_i}{\partial \eta_\mu}-\frac{\partial f^i_{r+1}}{\partial \chi_i}f_i=0.
	\end{equation}
	For each $i=1,\ldots,r$, the equation (\ref{Eq:PDEFro}) is a linear PDE. Due to the dependence of its coefficients, it can be considered to be defined on $\mathbb{R}^{s+1}$ and it always admits  a solution by the method of characteristics. Hence, we obtain that
	$$
	[f_iX_i,f_jX_j]=0,\qquad i,j=1,\ldots,r+1.
	$$
	By the induction hypothesis, our theorem follows.
\end{proof}

\end{document}
